\chapter{Bronchitis Treatment}

\section*{Definition and Overview}
Bronchitis is inflammation of the larger airways (the bronchi) that conduct air into the lungs. Treatment differs sharply between its two main forms. Acute bronchitis is a short-lived illness, usually caused by a viral infection, that settles on its own within a few weeks and generally requires only symptomatic care. Chronic bronchitis is a long-term component of chronic obstructive pulmonary disease (COPD), almost always related to smoking, and requires sustained management to control symptoms, reduce exacerbations, and slow progression. This chapter explains how each form is treated, the self-care measures that help, and the circumstances in which medical care is needed.

\section*{Epidemiology}
Acute bronchitis is extremely common and affects people of all ages, with most episodes occurring in autumn and winter when respiratory viruses circulate. It accounts for a large proportion of primary care consultations each year, and most adults experience an episode every few years. Chronic bronchitis is far more common in older adults, particularly those over 40, and is strongly linked to smoking, which causes the great majority of cases. It is somewhat more common in men in many populations, and long-term occupational exposure to dust, fumes, or air pollution increases the risk. Chronic bronchitis is a leading cause of chronic disability and hospital admission worldwide.

\section*{Aetiology and Pathophysiology}
Acute bronchitis is usually caused by viruses, including those responsible for colds and influenza, which infect and inflame the lining of the bronchi. Bacteria cause only a minority of cases and tend to follow a viral illness in people with chronic lung disease or impaired immunity. The inflammation makes the airway lining swell and produce excess mucus, narrowing the bronchi and producing cough and wheeze.

Chronic bronchitis develops after years of airway irritation, most often from tobacco smoke. Persistent inflammation leads to enlargement of the mucus-secreting glands, thickening of the airway walls, and gradual destruction of normal airway structure, so the airways become permanently narrowed and prone to infection. Contributing factors include air pollution, occupational dusts and fumes, and repeated respiratory infections.

\section*{Clinical Features}
The main symptom of both forms is a cough, frequently productive of phlegm.

\begin{itemize}
    \item \textbf{Acute bronchitis:} a dry or productive cough lasting one to three weeks, often accompanied by a mild fever, sore throat, runny nose, and chest discomfort
    \item \textbf{Chronic bronchitis:} a daily cough producing phlegm for at least three months of the year, in two consecutive years
    \item Shortness of breath, particularly on exertion
    \item Wheezing or a whistling sound when breathing out
    \item Fatigue and general malaise
    \item In chronic bronchitis, frequent chest infections and gradually worsening breathlessness
\end{itemize}

\section*{Differential Diagnosis}
A cough with breathlessness has many possible causes, so other conditions must be considered before bronchitis is confirmed.

\begin{itemize}
    \item \textbf{Pneumonia:} typically causes fever, breathlessness, and signs of lung consolidation
    \item \textbf{Asthma:} episodic wheezing triggered by allergens, exercise, or viral infections
    \item \textbf{Pertussis (whooping cough):} severe, paroxysmal coughing bouts that may end in a whoop or vomiting
    \item \textbf{Heart failure:} breathlessness and cough with frothy sputum, especially in older people with cardiac disease
    \item \textbf{Lung cancer:} a persistent or changing cough, weight loss, or haemoptysis, particularly in smokers
    \item \textbf{Tuberculosis and bronchiectasis:} chronic infections causing persistent cough and sputum
    \item \textbf{Gastro-oesophageal reflux:} acid reflux can drive a chronic, often nocturnal cough
\end{itemize}

\section*{When to Seek Medical Care}
See a doctor if a cough lasts more than three weeks, if blood is coughed up, if there is a high fever, if episodes of bronchitis recur, or if breathing is getting worse. People with asthma, COPD, or other long-term chest conditions should seek advice early when symptoms change.

\textbf{Call 000 (or local emergency services) for:}
\begin{itemize}
    \item Severe difficulty breathing or gasping for air
    \item Chest pain that is new, severe, or associated with breathlessness
    \item Coughing up a significant amount of blood
    \item Lips or face turning blue or grey
    \item Drowsiness, confusion, or extreme weakness
\end{itemize}

\section*{Diagnosis and Investigations}
Most cases of acute bronchitis are diagnosed from the history and examination alone; tests are arranged to exclude other causes or to assess chronic disease.

\begin{itemize}
    \item \textbf{History and examination:} the nature, duration, and triggers of the cough, smoking history, and chest auscultation
    \item \textbf{Oxygen saturation:} pulse oximetry to check blood oxygen levels
    \item \textbf{Chest X-ray:} to look for pneumonia or other lung disease when suspected
    \item \textbf{Spirometry:} breathing tests to diagnose and monitor chronic bronchitis or COPD
    \item \textbf{Sputum culture:} when a bacterial infection is suspected or the sputum is unusually coloured
    \item \textbf{Blood tests:} occasionally inflammatory markers or other investigations when the diagnosis is unclear
\end{itemize}

\section*{Management}
Treatment aims to relieve symptoms in acute disease and to control symptoms and slow progression in chronic disease.

\begin{itemize}
    \item \textbf{Self-care for acute bronchitis:} rest, adequate fluids, and simple analgesia such as paracetamol or ibuprofen for fever and discomfort
    \item \textbf{Cough remedies:} honey and warm drinks may help; cough medicines are of limited proven benefit and are best reserved for troublesome symptoms
    \item \textbf{Antibiotics:} not routinely given for acute bronchitis because most cases are viral; they are reserved for confirmed or strongly suspected bacterial infection or pneumonia
    \item \textbf{Smoking cessation:} the single most important treatment for chronic bronchitis, and the only intervention proven to slow disease progression
    \item \textbf{Inhalers:} reliever (bronchodilator) inhalers to open the airways, with preventer inhalers for persistent symptoms or frequent exacerbations
    \item \textbf{Pulmonary rehabilitation:} a structured programme of exercise, education, and breathing techniques that improves fitness and quality of life in chronic disease
    \item \textbf{Vaccination:} annual influenza and pneumococcal vaccination to reduce the risk of secondary infection
    \item \textbf{Follow-up:} regular review to monitor symptoms, lung function, and treatment response
\end{itemize}

\section*{Complications}
Acute bronchitis is usually self-limiting, but complications can occur.

\begin{itemize}
    \item Secondary bacterial infection leading to pneumonia
    \item Worsening of underlying asthma or COPD, sometimes requiring hospital care
    \item Respiratory failure in severe or advanced chronic disease
    \item Progression of chronic bronchitis to advanced COPD with persistent breathlessness and low oxygen levels
    \item Complications of repeated antibiotic or corticosteroid courses when these are prescribed frequently
    \item Reduced physical activity, deconditioning, and impaired quality of life
\end{itemize}

\section*{Prognosis and Outcomes}
Acute bronchitis almost always resolves fully within two to three weeks, although the cough may linger for several more weeks. Chronic bronchitis cannot be cured, but stopping smoking and using inhaled therapy can slow the decline in lung function and reduce symptoms. Many people with mild to moderate disease live active lives for many years. The outlook is considerably less favourable if smoking continues or if significant lung damage has already occurred.

\section*{Prevention and Self-Care}
The most effective prevention is avoiding the causes of airway irritation.

\begin{itemize}
    \item Do not smoke, and avoid breathing in other people's tobacco smoke
    \item Avoid exposure to dust, chemical fumes, and heavy air pollution where possible
    \item Practise good hand hygiene and keep away from others when unwell to reduce infection spread
    \item Keep vaccinations up to date, including influenza and pneumococcal vaccines if recommended
    \item During an episode, rest, drink warm fluids, and use steam or a humidifier if the air is dry
    \item Manage coexisting asthma or other chest conditions with a written action plan
    \item See a doctor for any cough lasting more than three weeks, or for new or worsening symptoms
\end{itemize}

\section*{Sources and Further Reading}
The following organisations provide reliable information on the treatment of bronchitis and related respiratory disease.

\begin{itemize}
    \item NHS. \textit{Bronchitis: treatment, self-care, and when to seek help.} https://www.nhs.uk/conditions/bronchitis/
    \item Mayo Clinic. \textit{Bronchitis: diagnosis and treatment.} https://www.mayoclinic.org/
    \item MedlinePlus. \textit{Bronchitis and cough treatment information.} https://medlineplus.gov/bronchitis.html
    \item MSD Manual. \textit{Acute bronchitis: treatment and professional guidance.} https://www.msdmanuals.com/
    \item Centers for Disease Control and Prevention. \textit{Antibiotic use and acute respiratory infections guidance.} https://www.cdc.gov/
    \item National Institute for Health and Care Excellence. \textit{Chronic obstructive pulmonary disease management guidance.} https://www.nice.org.uk/
\end{itemize}
